\documentclass{article}

\usepackage{microtype}
\usepackage{graphicx}
\usepackage{subcaption}
\usepackage{booktabs}
\usepackage{hyperref}
\usepackage{xspace}

\newcommand{\theHalgorithm}{\arabic{algorithm}}

% CN: ICML 2027 author kit 尚未发布；当前使用最新的官方 ICML 2026 review style。
% EN: The ICML 2027 author kit is not yet available; this draft uses the latest official ICML 2026 review style.
\usepackage{icml2026}

\usepackage{amsmath}
\usepackage{amssymb}
\usepackage{mathtools}
\usepackage{amsthm}

\newcommand{\ours}{\texttt{SHUTTLE}\xspace} 
\newcommand{\subsolver}{\texttt{PDUOT}\xspace}
\newcommand{\build}{\texttt{BuildHier}\xspace}
\newcommand{\init}{\texttt{Init}\xspace}
\newcommand{\pricing}{\texttt{Pricing}\xspace}
\newcommand{\prune}{\texttt{Prune}\xspace}
\newcommand{\pduotmap}{\ensuremath{T}\xspace}
\newcommand{\checkstop}{\texttt{CheckStop}\xspace}
\newcommand{\checkrestart}{\texttt{CheckRestart}\xspace}
\newcommand{\updateweight}{\texttt{UpdateWeight}\xspace}
\newcommand{\mm}{\texttt{MM}\xspace}
\newcommand{\hello}{\texttt{HELLO}\xspace}
\newcommand{\sinkhorn}{\texttt{Sinkhorn}\xspace}
\newcommand{\KL}{\operatorname{KL}}
\newcommand{\supp}{\operatorname{supp}}

\newtheorem{proposition}{Proposition}
\newtheorem{theorem}{Theorem}
\newtheorem{corollary}{Corollary}
\theoremstyle{definition}
\newtheorem{definition}{Definition}
\theoremstyle{plain}

\icmltitlerunning{\ours: Linear-Memory UOT via Edge Localization}

\begin{document}

\twocolumn[
%   \icmltitle{\ours: Sparse Hierarchical Unbalanced Optimal Transport Through Linear-Memory Edge Localization}
    \icmltitle{A Linear-Memory Hierarchical Solver for Unbalanced Optimal Transport}

    \begin{icmlauthorlist}
        \icmlauthor{Anonymous Authors}{anon}
    \end{icmlauthorlist}

    \icmlaffiliation{anon}{Anonymous Institution}
    \icmlcorrespondingauthor{Anonymous Author}{anon@example.com}
    \icmlkeywords{unbalanced optimal transport, large-scale optimization, hierarchical methods, sparse Sinkhorn}

    \vskip 0.3in
]

\printAffiliationsAndNotice{}

\begin{abstract}
% 问题背景和难点：说明 UOT 的作用以及无熵离散问题的规模瓶颈。
Unbalanced optimal transport (UOT) compares distributions while allowing their total mass to vary, yet its discrete formulation still involves quadratically many transport variables.
% 核心框架来源：将 balanced OT 的 dual-guided hierarchy 适配到 UOT。
We introduce \ours, which adapts the dual-guided hierarchical framework of balanced OT to large-scale unregularized UOT.
% 结构性依据与规模优势：共享的约束支持层间和层内定位边，并维持线性内存。
Exploiting the shared sparse-solution structure and dual constraints of balanced OT and UOT, \ours localizes edges across and within hierarchy levels while maintaining linear memory complexity.
% UOT 专门适配：校准传播的势，并开发 restricted UOT 子问题求解器。
We calibrate propagated potentials to account for soft marginal penalties and develop a tailored GPU-friendly restricted UOT solver for resulting subproblems.
% 理论保证：保留证明所需条件，但压缩为摘要级短语。
Under suitable conditions, we prove convergence of the subproblem solver and finite-tolerance termination of the hierarchical refinement.
% 实验结果：合并三层实验，突出效率和可扩展性。
Experiments on synthetic and real-world transport problems demonstrate accurate solutions, linear memory growth, and efficient scaling beyond dense solvers.
% 扩展能力：UFGW 与 sparse Sinkhorn 依赖 localized support，放在摘要结尾简洁说明。
We further extend the localized-support framework to sparse entropic refinement and UOT variants, obtaining efficient large-scale solvers with downstream behavior comparable to dense Sinkhorn on biological data.
\end{abstract}

\section{Introduction}

Optimal transport (OT) provides a geometric framework for comparing distributions and aligning datasets, while unbalanced optimal transport (UOT) relaxes marginal constraints to allow mass variation~\citep{chizat2018scaling,sejourne2023unbalanced}.
% CN: 先介绍单细胞应用，再列举其他领域及其需要边际放松的原因。
% EN: Introduce single-cell applications first, followed by other fields and their motivations for marginal relaxation.
This flexibility supports applications in single-cell genomics~\citep{schiebinger2019optimal,klein2025mapping}, domain adaptation~\citep{fatras2021unbalanced}, and point-cloud motion estimation~\citep{puy2020flot}, accommodating cell proliferation and death, mismatched class proportions, and occlusions.

Let $\mathbf{a}\in\mathbb R_{++}^m$ and $\mathbf{b}\in\mathbb R_{++}^n$ be the reference masses of two discrete measures, and let $\mathbf{C}\in\mathbb R^{m\times n}$ be their pairwise cost matrix.
UOT minimizes transport cost plus divergence penalties on the deviations of the transported marginals from $\mathbf{a}$ and $\mathbf{b}$.
We focus on KL marginal penalties, with other penalties discussed in Appendix~\ref{app:extension_penalty}:
\begin{equation}
\label{eq:primal-uot}
\min_{X\geq 0}\quad \langle \mathbf{C},\mathbf{X}\rangle+\rho_a\KL(\mathbf{X}\mathbf{1}_n\,|\,\mathbf{a})
+\rho_b\KL(\mathbf{X}^\top\mathbf{1}_m\,|\,\mathbf{b}),
\end{equation}
Here, $\mathbf{X}\in\mathbb R_+^{m\times n}$ is the transport plan, $\rho_a,\rho_b>0$ control the marginal penalties, and $\KL(\mathbf{p}\,|\,\mathbf{h})=\sum_k(p_k\log(p_k/\mathbf{h}_k)-p_k+\mathbf{h}_k)$.

The dual of~\eqref{eq:primal-uot} optimizes potentials $\mathbf{f}\in\mathbb R^m$ and $\mathbf{g}\in\mathbb R^n$:
\begin{equation} 
\label{eq:dual-uot}
\begin{aligned}
&D(\mathbf{f},\mathbf{g}):=\rho_a\langle \mathbf{a},\mathbf{1}_m-e^{-\mathbf{f}/\rho_a}\rangle
+\rho_b\langle \mathbf{b},\mathbf{1}_n-e^{-\mathbf{g}/\rho_b}\rangle,\\
&\max_{\mathbf{f},\mathbf{g}}D(\mathbf{f},\mathbf{g})\quad\text{s.t.}\quad f_i+g_j\leq C_{ij}\quad\forall i,j.
\end{aligned}
\end{equation}
The primal formulation has one transport variable for each of the $mn$ source--target pairs; the dual has only $m+n$ potentials but still imposes $mn$ constraints.
This quadratic problem size poses substantial memory and computational challenges at large sample sizes.

\paragraph{Related work.}
Regularization methods add penalties to the transport plan to facilitate optimization, including squared-$\ell_2$ regularization~\citep{nguyen2023unbalanced} and, most prominently, entropic regularization.
The latter enables generalized \sinkhorn iterations, including stabilized and translation-invariant scaling variants~\citep{chizat2018scaling,schmitzer2019stabilized,pham2020unbalanced,sejourne2022faster}.
However, additional regularization introduces bias relative to unregularized UOT, and entropic regularization generally yields a dense coupling whose explicit materialization requires $\Theta(mn)$ memory.

Low-rank methods constrain the coupling to admit a compact nonnegative factorization, improving computational and memory scalability at a fixed rank~\citep{scetbon2023unbalanced,halmos2024low}.
This scalability comes at the cost of potentially substantial approximation error when the chosen rank is insufficient to represent the optimal coupling.

A third line of work solves the original UOT problem.
Majorization--minimization (\mm) directly updates the coupling through multiplicative iterations~\citep{chapel2021unbalanced}.
Alternatively, balanced-OT reductions separate marginal estimation from coupling construction through upfront estimation~\citep{liu2026solving} or iterative Frank--Wolfe refinement~\citep{sejourne2022faster,xia2026dual}.
Another approach exploits coupling sparsity through a network-simplex framework that maintains forest-supported transport plans~\citep{agazzotti2026network}.
Scalability remains limited by quadratic coupling storage in standard \mm implementations, the cost of marginal estimation or sequential balanced OT subproblems in reduction methods, and sequential tree updates in network simplex that hinder GPU parallelization.
This leaves open how to solve the large scale unregularized UOT problem with linear memory and efficient GPU parallelism.

\paragraph{Our approach.}
We introduce \ours, a \textbf{S}parse \textbf{H}ierarchical solver for large-scale \textbf{U}nbalanced Optimal \textbf{T}ransport \textbf{T}hrough \textbf{L}inear-Memory \textbf{E}dge Localization.
UOT admits sparse optimal couplings and, for convex marginal penalties, shares the dual constraints $f_i+g_j \leq C_{ij}$ of balanced OT, which motivates extending the dual-guided hierarchy of \hello to the unbalanced setting~\citep{xia2026dual}.
The marginal penalties, however, turn the restricted transport linear programs into nonlinear convex problems.
We therefore develop \subsolver, a GPU-friendly primal--dual solver for restricted UOT.
Combining \subsolver with hierarchical edge localization and calibration of propagated dual potentials yields a GPU-friendly UOT solver with linear memory.

Our contributions are summarized as follows:
\begin{itemize}
    \item \textit{Dual-guided hierarchy for UOT:} We extend dual-guided hierarchical edge localization to unregularized UOT, enabling large-scale optimization with linear memory and finite-tolerance termination guarantees for safeguarded refinement.
    \item \textit{GPU-efficient restricted solver:} We develop \subsolver, an efficient GPU-parallel primal--dual solver for restricted UOT, with convergence guarantees for its fixed-metric core.
    \item \textit{Extensions and applications:} We extend \ours to efficiently approximate entropic UOT at small regularization strengths and solve unbalanced fused Gromov--Wasserstein (UFGW) problems, demonstrating their effectiveness on large-scale biological alignment tasks.
\end{itemize}

\section{Sparse Hierarchical Solver for UOT}
\label{sec:method}


\subsection{Overall Framework}
\label{sec:framework}

Let $\mathcal S=\{\mathbf{s}_i\}_{i=1}^m$ and $\mathcal D=\{\mathbf{d}_j\}_{j=1}^n$ denote the source and target supports in~\eqref{eq:primal-uot}, with $C_{ij}=c(\mathbf{s}_i,\mathbf{d}_j)$.
With the transported marginals $\mathbf{X}\mathbf{1}_n$ and $\mathbf{X}^\top\mathbf{1}_m$ fixed, the marginal penalties are constant and minimizing over $\mathbf{X}$ reduces to balanced OT.
Consequently, whenever an optimum exists, there is an optimal coupling with at most $m+n-1$ positive entries.
This sparsity motivates optimizing over an active support $\mathcal{N}\subseteq[m]\times[n]$, with all other transport variables fixed to zero.
The central task is to discover and refine this support without maintaining all $mn$ variables.

The dual formulation~\eqref{eq:dual-uot} provides a pricing mechanism for this task.
Each transport variable $X_{ij}$ corresponds to a linear dual inequality $f_i+g_j\leq C_{ij}$, so adding a primal variable amounts to enforcing an additional dual constraint.
This correspondence allows us to extend the restricted-master and pricing strategy of LP column generation to UOT, with a nonlinear convex restricted master.
To our knowledge, this form of column generation remains comparatively unexplored for unregularized KL-UOT.

The KKT conditions for~\eqref{eq:primal-uot} and~\eqref{eq:dual-uot} are
\begin{equation*}
\begin{aligned}
\mathbf{X}\mathbf{1}_n&=\mathbf{a}\odot e^{-\mathbf{f}/\rho_a},\, \mathbf{X}^\top\mathbf{1}_m=\mathbf{b}\odot e^{-\mathbf{g}/\rho_b},\\
f_i+g_j&\leq C_{ij},\\
X_{ij}(f_i+g_j-C_{ij})&=0.
\end{aligned}
\end{equation*}

An exact restricted solution satisfies these conditions on $\mathcal{N}$, together with the marginal-response equations on all nodes.
It is therefore globally optimal whenever its potentials also satisfy every omitted dual inequality.
We use the dual score
\begin{equation}
\label{eq:dual-score}
\sigma_{ij}:=f_i+g_j-C_{ij}
\end{equation}
to price omitted edges: a positive score identifies a violated constraint and a candidate variable for the next restricted solve.
The resulting refinement alternates restricted UOT solves with full-space pricing and edge insertion, while budgeted pruning controls the active-support size.

We organize this refinement in the dual-guided hierarchy introduced by \hello~\citep{xia2026dual}.
The finest problem is indexed by $0$, and the coarsest problem is indexed by $L$.
Adjacent levels change one marginal at a time, allowing the unchanged potential to be inherited directly.
At each finer level, propagated and calibrated dual potentials guide the initial active support, which is then refined using the same dual scores.
We write $\mathcal Z^{(\ell)}=(\mathbf{X}^{(\ell)},\mathbf{f}^{(\ell)},\mathbf{g}^{(\ell)})$ for a level's primal--dual state and $\mathcal E^{(\ell)}$ for its complete edge set.
Within a level, $\mathcal N$, $\mathcal B$, and $\mathcal Z$ denote the current active support, structural edges, and primal--dual state.
Algorithm~\ref{alg:ours} summarizes the framework.

\begin{algorithm}[t]
\caption{\ours}
\label{alg:ours}
\begin{algorithmic}[1]
\STATE \textbf{Input:} KL-UOT problem $\mathcal P^{(0)}$, support budgets $\kappa,\gamma,\beta$, and tolerance $\varepsilon$.
\STATE \textbf{Output:} finest-level primal--dual state $\mathcal Z^{(0)}$.
\medskip
\STATE $\{\mathcal P^{(\ell)}\}_{\ell=0}^{L}\gets\build(\mathcal P^{(0)})$
\STATE $\mathcal Z^{(L)}\gets\subsolver(\mathcal P^{(L)},\mathcal E^{(L)},\varnothing;\varepsilon)$
\FOR{$\ell=L-1,\ldots,0$}
    \STATE $(\mathcal N,\mathcal B,\mathcal Z)\gets\init(\mathcal P^{(\ell)},\mathcal Z^{(\ell+1)};\kappa)$
    \LOOP
        \STATE $\mathcal Z\gets\subsolver(\mathcal P^{(\ell)},\mathcal N,\mathcal Z;\varepsilon)$
        \STATE $(\mathcal V,\eta)\gets\pricing(\mathcal P^{(\ell)},\mathcal Z;\gamma)$
        \IF{$\eta\leq\varepsilon$}
            \STATE \textbf{break}
        \ENDIF
        \STATE $\mathcal N\gets\mathcal N\cup\mathcal V$
        \STATE $\mathcal N\gets\prune(\mathcal N,\mathcal Z,\mathcal B,\mathcal V;\beta)$
    \ENDLOOP
    \STATE $\mathcal Z^{(\ell)}\gets\mathcal Z$
\ENDFOR
\medskip
\STATE \textbf{Return:} $\mathcal Z^{(0)}$.
\end{algorithmic}
\end{algorithm}

\subsection{Restricted UOT with \subsolver}
\label{sec:pduot}

Algorithm~\ref{alg:ours} repeatedly solves UOT on the current active support, giving rise to the following restricted problem.
For columnwise vectorization of $\mathbf{X}$, the marginal-sum matrix is
\begin{equation}
\label{eq:transport-incidence}
\mathbf{A}=\begin{pmatrix}\mathbf{1}_n^\top\otimes \mathbf{I}_m\\\mathbf{I}_n\otimes\mathbf{1}_m^\top\end{pmatrix}\in\mathbb R^{(m+n)\times mn}.
\end{equation}
Thus $\mathbf{A}\operatorname{vec}(\mathbf{X})=(\mathbf{X}\mathbf{1}_n,\mathbf{X}^\top\mathbf{1}_m)$, with the two vectors stacked vertically.
Let $\mathbf{A}_{\mathcal N}$ retain the columns indexed by $\mathcal N$, and let $\mathbf{c}_{\mathcal N}$ retain the corresponding entries of $\operatorname{vec}(\mathbf{C})$.
Define $G(\mathbf{r},\mathbf{q})=\rho_a\KL(\mathbf{r}\,|\,\mathbf{a})+\rho_b\KL(\mathbf{q}\,|\,\mathbf{b})$ for $\mathbf{r}\in\mathbb R^m$ and $\mathbf{q}\in\mathbb R^n$.

\begin{definition}[Restricted UOT]
\label{def:restricted-uot}
For an active support $\mathcal N\subseteq[m]\times[n]$, restricted UOT is~\eqref{eq:primal-uot} with $X_{ij}=0$ for every $(i,j)\notin\mathcal N$.
Let $\mathbf{x}\in\mathbb R^{|\mathcal N|}$ collect the transport variables indexed by $\mathcal N$; the primal and dual problems are
\begin{align}
\label{eq:restricted-uot}
&\min_{\mathbf{x}\geq0}\;\langle \mathbf{c}_{\mathcal{N}},\mathbf{x}\rangle+G(\mathbf{A}_{\mathcal N}\mathbf{x}),\\
\label{eq:restricted-dual-uot}
&\max_{\mathbf{f},\mathbf{g}}\;D(\mathbf{f},\mathbf{g})\quad\text{s.t. }f_i+g_j\leq C_{ij},\ (i,j)\in\mathcal N.
\end{align}
The dual objective is unchanged from~\eqref{eq:dual-uot}, while only the inequalities indexed by $\mathcal N$ are retained.
\end{definition}

These repeated solves must efficiently produce high-accuracy primal--dual solutions satisfying the unregularized restricted KKT conditions.
Adapting existing UOT algorithms to active supports presents challenges in jointly achieving this accuracy and computational efficiency (see Section~\ref{sec:restricted-experiments}).
We therefore develop \subsolver, a practical GPU-efficient restricted UOT solver based on primal--dual hybrid gradient (PDHG)~\citep{esser2010general,chambolle2011first}.
The interface $\subsolver(\mathcal P,\mathcal N,\mathcal Z_{\mathrm{init}};\varepsilon)$ accepts an initial primal--dual state and returns a state satisfying the restricted stopping criterion; $\varnothing$ selects the default initialization.

The linear transport cost and separable marginal penalties lead to a saddle formulation with simple primal and dual updates.
With
\begin{equation*}
F(\mathbf{x})=\langle \mathbf{c}_{\mathcal{N}},\mathbf{x}\rangle+\iota_{\mathbb R_+^{|\mathcal{N}|}}(\mathbf{x}),
\end{equation*}
the saddle formulation is
\begin{equation}
\label{eq:restricted-saddle}
\min_{\mathbf{x}}\max_{\mathbf{y}}\;F(\mathbf{x})+\langle \mathbf{A}_{\mathcal N}\mathbf{x},\mathbf{y}\rangle-G^*(\mathbf{y}),
\end{equation}
where $\mathbf{y}=-(\mathbf{f},\mathbf{g})$ and
\begin{equation*}
G^*(\mathbf{y})=\rho_a\sum_i a_i(e^{y_i/\rho_a}-1)+\rho_b\sum_j b_j(e^{y_{m+j}/\rho_b}-1).
\end{equation*}

Let $\mathbf{D}_x$ and $\mathbf{D}_y$ be positive diagonal step matrices.
\subsolver uses the primal-first PDHG map $\pduotmap(\mathbf{x},\mathbf{y})=(\widehat{\mathbf{x}},\widehat{\mathbf{y}})$:
\begin{equation}
\label{eq:pduot-map}
\begin{aligned}
\widehat{\mathbf{x}}&=\left[\mathbf{x}-\mathbf{D}_x(\mathbf{c}_{\mathcal{N}}+\mathbf{A}_{\mathcal N}^\top \mathbf{y})\right]_+,\\
\widehat{\mathbf{y}}&=\operatorname{prox}_{G^*}^{\mathbf{D}_y^{-1}}\left(\mathbf{y}+\mathbf{D}_y\mathbf{A}_{\mathcal N}(2\widehat{\mathbf{x}}-\mathbf{x})\right).
\end{aligned}
\end{equation} 
Here, $\operatorname{prox}_{G^*}^{\mathbf{D}_y^{-1}}(\mathbf{u})=\arg\min_v\{G^*(\mathbf{v})+\tfrac12\|\mathbf{v}-\mathbf{u}\|_{\mathbf{D}_y^{-1}}^2\}$.
The primal step is a nonnegative projection driven by the dual scores, while the dual proximal map separates over nodes.
For an input $u$, reference mass $h>0$, penalty $\rho$, and diagonal dual step $\mathbf{s}>0$, each coordinate solves
\begin{equation*}
v-u+\mathbf{s} h e^{v/\rho}=0.
\end{equation*}
We solve this equation in the log domain to avoid overflow; Appendix~\ref{app:pduot-details} gives the scalar formulation.

We combine this PDHG map with the following practical components, following restarted primal--dual methods~\citep{lu2407restarted,lin2025practical}.
\begin{itemize}
\item \textit{Reflected Halpern iteration.}
Within each restart epoch, we apply
\begin{equation}
\label{eq:halpern}
\mathbf{z}^{k+1}=\frac{1}{k+2}\mathbf{z}^{\mathrm a}+\frac{k+1}{k+2}(2\pduotmap \mathbf{z}^k-\mathbf{z}^k),
\end{equation}
where $\mathbf{z}=(\mathbf{x},\mathbf{y})$, $\mathbf{z}^{\mathrm a}$ is the fixed epoch anchor, and $k$ counts iterations within the epoch.
We use the mapped point $\widehat{\mathbf{z}}=\pduotmap \mathbf{z}^k$ for stopping and output, since its coupling is nonnegative.

\item \textit{Diagonal preconditioning.}
We set $\mathbf{D}_x=w\overline{\mathbf{D}}_x$ and $\mathbf{D}_y=\overline{\mathbf{D}}_y/w$, where $w>0$ balances primal and dual steps.
For node-covering supports, the base steps adapt to the active node degrees $d_v$:
\begin{equation}
\label{eq:pduot-metric}
(\overline{\mathbf{D}}_x)_{ee}=\frac{\zeta}{2},\qquad
(\overline{\mathbf{D}}_y)_{vv}=\frac{\zeta}{d_v},\qquad 0<\zeta<1.
\end{equation}
This choice ensures $\|\mathbf{D}_y^{1/2}\mathbf{A}_{\mathcal N}\mathbf{D}_x^{1/2}\|_2<1$.

\item \textit{Adaptive restart.}
We monitor the fixed-point residual and restart after sufficient reduction, moderate reduction followed by deterioration, or an epoch-length threshold.
Each restart resets the anchor to the mapped point and the epoch counter to zero.

\item \textit{Primal-weight adjustment.}
At restart boundaries, a bounded log-domain proportional--integral controller adjusts $w$ using primal and dual displacements measured in the base metrics.
The base preconditioners remain fixed throughout the solve.
\end{itemize}
Algorithm~\ref{alg:pduot} summarizes the solver; Appendix~\ref{app:pduot-details} gives the practical control rules.
At the mapped point, \checkstop evaluates marginal consistency $\delta_{\mathrm m}$, stationarity $\delta_{\mathrm s}$, and relative gap $\delta_{\mathrm g}$ in the original problem coordinates, returning when
\begin{equation}
\label{eq:pduot-stopping}
\max\{\delta_{\mathrm m},\delta_{\mathrm s},\delta_{\mathrm g}\}\leq\varepsilon.
\end{equation}
Appendix~\ref{app:stopping-criterion} defines these measures and explains their relation to the restricted KKT conditions and the balanced-OT limit.

\begin{algorithm}[t]
\caption{\subsolver for restricted UOT}
\label{alg:pduot}
\begin{algorithmic}[1]
\STATE \textbf{Input:} UOT problem $\mathcal P$, active support $\mathcal N$, initial state $\mathbf{z}^0$, tolerance $\varepsilon$.
\STATE \textbf{Output:} restricted primal--dual solution $(\widehat{\mathbf{x}},\widehat{\mathbf{f}},\widehat{\mathbf{g}})$.
\STATE Choose $\overline{\mathbf{D}}_x,\overline{\mathbf{D}}_y,w$; set $\mathbf{z}^{\mathrm a}\gets \mathbf{z}^0$.
\LOOP
    \STATE $\mathbf{D}_x\gets w\overline{\mathbf{D}}_x$, $\mathbf{D}_y\gets\overline{\mathbf{D}}_y/w$
    \STATE $\mathbf{z}\gets \mathbf{z}^{\mathrm a}$, $k\gets0$
    \LOOP
        \STATE $\widehat{\mathbf{z}}\gets\pduotmap \mathbf{z}$ by~\eqref{eq:pduot-map}; retain $\boldsymbol{\lambda}$.
        \IF{$\checkstop(\widehat{\mathbf{z}},\boldsymbol{\lambda};\varepsilon)$}
            \STATE $(\widehat{\mathbf{f}},\widehat{\mathbf{g}})\gets-\widehat{\mathbf{y}}$
            \STATE \textbf{Return:} $(\widehat{\mathbf{x}},\widehat{\mathbf{f}},\widehat{\mathbf{g}})$.
        \ENDIF
        \IF{$\checkrestart(\mathbf{z},\widehat{\mathbf{z}},\mathbf{z}^{\mathrm a},k)$}
            \STATE \textbf{break}
        \ENDIF
        \STATE $\mathbf{z}\gets\dfrac{\mathbf{z}^{\mathrm a}}{k+2}+\dfrac{k+1}{k+2}(2\widehat{\mathbf{z}}-\mathbf{z})$
        \STATE $k\gets k+1$
    \ENDLOOP
    \STATE $w\gets\updateweight(\widehat{\mathbf{z}},\mathbf{z}^{\mathrm a},w)$
    \STATE $\mathbf{z}^{\mathrm a}\gets\widehat{\mathbf{z}}$
\ENDLOOP
\end{algorithmic}
\end{algorithm}

The following proposition gives convergence guarantees for the exact, fixed-metric core of \subsolver.

\begin{proposition}[Convergence of \subsolver]
\label{prop:pduot-convergence}
Fix an active support covering all nodes, and suppose $\mathbf a,\mathbf b>0$, $\rho_a,\rho_b>0$, and all active costs are finite.
Let the positive diagonal steps be fixed and satisfy $\|\mathbf D_y^{1/2}\mathbf A_{\mathcal N}\mathbf D_x^{1/2}\|_2<1$, and assume exact proximal evaluations.
Then the following statements hold:
\begin{itemize}
\item \textbf{Global convergence.}
The unrestarted iteration~\eqref{eq:halpern} and its mapped points converge to a restricted KKT point from any initial anchor.
\item \textbf{Local linear convergence.}
Initialized sufficiently near any restricted KKT point, the restart anchors converge linearly to the restricted KKT solution set under either a sufficiently long fixed restart period or the residual-reduction rule in Appendix~\ref{app:pduot-restart}.
Each restart sets the new anchor to the mapped point.
\end{itemize}
\end{proposition}

Appendices~\ref{app:pduot-global}--\ref{app:pduot-restart} give the proof and precise restart rules.
Practical control steps are detailed in Appendix~\ref{app:pduot-controls}.

\subsection{UOT-Aware Hierarchy and Refinement}
\label{sec:uot-hierarchy}

\paragraph{Potential propagation and calibration.}
\build follows the nested construction of \hello: one fixed permutation is applied to each marginal, and each coarser level truncates the currently larger side~\citep{xia2026dual}.
\init combines the potential propagation and calibration below with initial active-support construction.
Suppose the target side is shared by two adjacent levels.
We inherit its potential $\mathbf{g}$ and complete the finer source potential by the $\mathbf{c}$-transform
\begin{equation*}
f_i=\min_j\{C_{ij}-g_j\}.
\end{equation*}
This completion is feasible for every pairwise dual constraint at the finer level.

The transformation $\mathbf{f}\leftarrow \mathbf{f}+t\mathbf{1}_m$ and $\mathbf{g}\leftarrow \mathbf{g}-t\mathbf{1}_n$ preserves all dual scores, while the KL-UOT dual objective varies with $t$.
Define
\begin{equation*}
Z_a=\sum_i a_i e^{-f_i/\rho_a},\qquad Z_b=\sum_j b_j e^{-g_j/\rho_b}.
\end{equation*}

\begin{proposition}[Translation calibration]
\label{prop:translation-calibration}
Along the score-preserving translation direction, the KL-UOT dual objective is strictly concave and is maximized at
\begin{equation*}
t^\star=\frac{\rho_a\rho_b}{\rho_a+\rho_b}\left(\log Z_a-\log Z_b\right).
\end{equation*}
At the calibrated potentials, the total transported masses implied by the two marginal responses are equal.
\end{proposition}

The result follows by setting the directional derivative $e^{-t/\rho_a}Z_a-e^{t/\rho_b}Z_b$ to zero.
Calibration retains every edge ranking while removing the global response mismatch introduced by cross-level propagation.

\paragraph{Initial active support.}
For every source node, we retain the $\kappa$ target nodes with largest propagated dual scores and apply the symmetric selection for every target node.
Their union gives the bidirectional dual assignment and covers both sides of the bipartite graph.
We additionally include structural edges $\mathcal B$ given by a northwest-corner skeleton constructed from the reference masses.
In KL-UOT, these structural edges provide persistent node coverage and improve the conditioning of repeated restricted solves; they do not enforce hard marginal equalities.
\init returns the union $\mathcal N$ of dual-assigned and structural edges, the protected set $\mathcal B$, and an initial state $\mathcal Z$ with zero flows and the propagated, calibrated potentials.

\paragraph{Refinement and stopping.}
At a fixed hierarchy level, \subsolver first solves~\eqref{eq:restricted-uot} on the current active support.
\pricing evaluates~\eqref{eq:dual-score} in tiles without materializing $\mathbf{C}$ or the dense coupling.
It returns a candidate set $\mathcal V$, formed by retaining up to $\gamma$ largest positive dual scores per source and target node, together with the full-space dual-violation residual $\eta$.
After inserting $\mathcal V$, \prune protects structural edges $\mathcal B$, numerically positive-flow edges in $\mathcal Z$, and the new candidates $\mathcal V$; remaining old edges are removed in increasing order of relevance to dual tightness under budget factor $\beta$.
Flows on retained edges and node potentials are warm-started in the next restricted solve, and new edges receive zero initial flow.

The restricted stopping criterion combines the marginal-response and stationarity residuals in~\eqref{eq:pduot-residuals} with the gap in~\eqref{eq:pduot-split-gap}.
Both the restricted solve and the full-space dual-violation test use the same tolerance $\varepsilon$.
Since \subsolver returns only after the restricted criterion is satisfied, the outer loop checks only $\eta\leq\varepsilon$.
The full-space scan also produces the candidates for the next round, so stopping and support insertion share one streamed pass.

\subsection{Convergence and Complexity}
\label{sec:analysis}

We first connect a restricted \subsolver solution to the full KL-UOT problem.

\begin{proposition}[Restricted-to-full recovery]
\label{prop:restricted-recovery}
Let $(\widehat{\mathbf{x}},\widehat{\mathbf{f}},\widehat{\mathbf{g}})$ satisfy the KKT conditions of~\eqref{eq:restricted-uot} on $\mathcal{N}$.
If $\widehat f_i+\widehat g_j\leq C_{ij}$ for every $(i,j)\in[m]\times[n]$, then extending $\widehat{\mathbf{x}}$ by zero outside $\mathcal{N}$ yields a global primal solution of~\eqref{eq:primal-uot}, and $(\widehat{\mathbf{f}},\widehat{\mathbf{g}})$ is a global dual solution of~\eqref{eq:dual-uot}.
\end{proposition}

Restricted KKT conditions already provide the marginal responses, active-edge feasibility, and complementarity.
The full-space inequalities supply the only missing KKT conditions.

Finite-tolerance refinement requires accounting for inexact restricted solves.
For a coupling $\mathbf{X}$ with positive transported marginals, define primal-consistent potentials and scores by
\begin{equation*}
\begin{aligned}
f_i^\mathbf{X}&=-\rho_a\log(r_i/a_i),\quad g_j^\mathbf{X}=-\rho_b\log(q_j/b_j),\\
\sigma_{ij}^\mathbf{X}&=f_i^\mathbf{X}+g_j^\mathbf{X}-C_{ij}.
\end{aligned}
\end{equation*}
The coordinate derivative satisfies $\nabla_{ij}\Phi(\mathbf{X})=-\sigma_{ij}^\mathbf{X}$.

\begin{proposition}[One-edge descent]
\label{prop:one-edge-descent}
Let $h_{ij}(\mathbf{X})=\rho_a/r_i+\rho_b/q_j$.
If $\sigma_{ij}^\mathbf{X}>0$, then $t=\sigma_{ij}^\mathbf{X}/h_{ij}(\mathbf{X})$ satisfies
\begin{equation*}
\Phi(\mathbf{X}+t\mathbf{E}_{ij})\leq\Phi(\mathbf{X})-\frac{(\sigma_{ij}^\mathbf{X})^2}{2h_{ij}(\mathbf{X})}.
\end{equation*}
\end{proposition}

The proposition follows from a second-order upper bound along the nonnegative coordinate direction $\mathbf{E}_{ij}$.
This direction is always primal feasible because KL-UOT penalizes its marginals rather than imposing hard marginal equalities.

The following theorem states a finite-tolerance guarantee for a safeguarded refinement variant.
At round $k$, let an inexact restricted oracle return $\mathbf{X}^k$ on $\mathcal{N}_k$ with restricted projected-gradient residual at most $\varepsilon$.
If the largest inactive primal-consistent score exceeds $\varepsilon$, the update selects such an edge $e_k$, forms the trial point from Proposition~\ref{prop:one-edge-descent}, and retains $\supp(\mathbf{X}^k)\cup\{e_k\}$ in the next support.

\begin{theorem}[Finite-tolerance refinement]
\label{thm:refinement}
Assume that $\Phi$ is bounded below by $\Phi_{\inf}$, $h_{e_k}(\mathbf{X}^k)\leq H$, and the next restricted solve returns $\mathbf{X}^{k+1}$ satisfying
\begin{equation*}
\begin{aligned}
\Phi(\mathbf{X}^{k+1})&\leq\Phi(\mathbf{X}^k+t_k\mathbf{E}_{e_k})+\xi_k,\\
0\leq\xi_k&\leq\theta\frac{(\sigma_{e_k}^{\mathbf{X}^k})^2}{2h_{e_k}(\mathbf{X}^k)}
\end{aligned}
\end{equation*}
for constants $H<\infty$ and $\theta\in[0,1)$.
Then the safeguarded refinement terminates after at most
\begin{equation}
\label{eq:refinement-bound}
N_{\mathrm{ref}}\leq\left\lceil\frac{2H(\Phi(\mathbf{X}^0)-\Phi_{\inf})}{(1-\theta)\varepsilon^2}\right\rceil
\end{equation}
nonterminal rounds and returns a coupling whose full projected-gradient residual is at most $\varepsilon$.
\end{theorem}

Each nonterminal round decreases the objective by at least $(1-\theta)\varepsilon^2/(2H)$, and telescoping proves~\eqref{eq:refinement-bound}.
The theorem isolates an implementable safeguard for inexact \subsolver calls.
The practical variant uses \subsolver potentials for candidate ranking; primal-consistent scores or two-sided response corrections can be used for the safeguarded stopping test.
The supplementary material will provide further details and explicit marginal bounds for $H$.

Finally, let $\mathbf{s}=|\mathcal{N}|$ denote the largest maintained active support.
\subsolver stores edge flows and sparse incidence indices in $O(\mathbf{s})$ memory, while the hierarchy stores input features and node quantities in $O((m+n)d+m+n)$ memory.
The tiled full-space scan uses the same persistent memory and performs $O(mn)$ arithmetic per pass.
Consequently,
\begin{equation}
\label{eq:memory-complexity}
\operatorname{Mem}(\ours)=O\bigl((m+n)d+m+n+\mathbf{s}\bigr).
\end{equation}
Under the configured support budget $\mathbf{s}\leq\beta(m+n)$ and fixed feature dimension,~\eqref{eq:memory-complexity} is linear in the number of input points.

\section{Extensions}
\label{sec:extensions}

The sparse support and restricted UOT solver provided by \ours also support entropic refinement and transport between relational datasets.

\subsection{Sparse Entropic Refinement}
\label{sec:entropic-extension}

Let $\Phi$ denote the unregularized objective in~\eqref{eq:primal-uot}, and let $\mathcal{N}_\star$ be the final active support returned by \ours.
To approximate entropic UOT at a small regularization strength $\varepsilon>0$, we solve
\begin{equation}
\label{eq:sparse-entropic-extension}
\min_{X\geq0,\,\supp(X)\subseteq\mathcal{N}_\star}\;\Phi(\mathbf{X})+\varepsilon\KL(\mathbf{X}\,|\,\mathbf{a}\otimes \mathbf{b}).
\end{equation}
Generalized Sinkhorn updates on this fixed support cost $O(|\mathcal{N}_\star|+m+n)$ time and memory per iteration.
Reusing the final active support avoids an additional support-construction scan and retains an explicitly sparse coupling for downstream tasks.
The approximation comes from restricting the support of the full entropic problem; a small $\varepsilon$ motivates reusing edges localized by the unregularized solver, while its downstream accuracy is evaluated experimentally.
Appendix~\ref{app:entropic-extension} describes the sparse updates and an optional support expansion.

\subsection{Unbalanced Fused Gromov--Wasserstein}
\label{sec:ufgw-extension}

UFGW combines feature matching with alignment of within-dataset relations.
Let $\mathbf{A}\in\mathbb R^{m\times m}$ and $\mathbf{B}\in\mathbb R^{n\times n}$ encode these relations, and let $\mathbf{M}\in\mathbb R^{m\times n}$ be the feature cost.
We use a two-coupling block-coordinate relaxation with tensorized KL marginal penalties~\citep{sejourne2023unbalanced}.
% CN: 请作者补入原始 UGW relaxation 论文的引用；此处暂用现有综述引用，不修改 bibliography。
% EN: Add the original UGW relaxation citation; the existing survey is cited provisionally without changing the bibliography.
For nonnegative couplings $\mathbf{P},\mathbf{Q}$, define
\begin{equation*}
\mathcal Q_{A,B}(\mathbf{P},\mathbf{Q})=\sum_{i,k,j,l}(A_{ik}-B_{jl})^2P_{ij}Q_{kl}
\end{equation*}
and $D_a(\mathbf{P},\mathbf{Q})=\KL((\mathbf{P}\mathbf{1})\otimes(\mathbf{Q}\mathbf{1})\,|\,\mathbf{a}\otimes \mathbf{a})$, with $D_b$ defined analogously for target marginals.
The relaxed objective is
\begin{equation}
\label{eq:ufgw-relaxation}
\begin{aligned}
\min_{P,Q\geq0}\quad &(1-\alpha)\mathcal Q_{A,B}(\mathbf{P},\mathbf{Q})\\
&+\frac{\alpha}{2}\langle \mathbf{M},\mathbf{P}+\mathbf{Q}\rangle
+\rho_aD_a(\mathbf{P},\mathbf{Q})+\rho_bD_b(\mathbf{P},\mathbf{Q}),
\end{aligned}
\end{equation}
where $\alpha\in[0,1]$ controls feature matching and $\alpha=0$ gives the UGW relaxation.
Fixing either coupling yields a KL-UOT subproblem whose marginal penalty coefficients are scaled by the mass of the fixed coupling.
We alternate \ours solves for the two blocks, with mass-balancing rescaling after each block update; sparse entropic refinement provides an optional regularized variant.
For the squared-Euclidean structural and feature costs used in our experiments, each block cost admits an exact squared-Euclidean embedding computed from sparse coupling moments.
This representation allows \ours to scan the block costs without storing dense structural or cross-cost matrices.
The block formulas, rescaling, and entropic variant are given in Appendix~\ref{app:ufgw-extension}.

\section{Experiments}
\label{sec:experiments}

\subsection{Restricted UOT Solvers}
\label{sec:restricted-experiments}

We compare \subsolver with existing solvers on shared restricted UOT instances, evaluating runtime and primal--dual accuracy across support sizes and marginal penalty parameters.
% CN: 待补充统一评测指标、基线配置以及运行时间与精度的对比结果。
% EN: Add shared evaluation metrics, baseline configurations, and runtime--accuracy comparisons.

\bibliography{references}
\bibliographystyle{icml2026}

\newpage
\appendix
\onecolumn
\section{PDUOT Analysis and Control Steps}
\label{app:pduot-details}

\paragraph{KL proximal evaluation.}
For dual input $u$, reference mass $h>0$, penalty $\rho>0$, and step $\mathbf{s}>0$, the coordinate proximal problem is
\begin{equation*}
\min_v\;\rho h(e^{v/\rho}-1)+\frac{(v-u)^2}{2\mathbf{s}}.
\end{equation*}
Its optimality equation $v-u+\mathbf{s} h e^{v/\rho}=0$ has the unique solution $v=u-\rho W((\mathbf{s} h/\rho)e^{u/\rho})$, where $W$ is the principal Lambert--$W$ function.
To avoid forming the exponential, let $\ell=\log(\mathbf{s} h/\rho)+u/\rho$ and solve $e^q+q=\ell$ for $q$.
Then $e^q=W(e^\ell)$ and $v=u-\rho e^q$; equivalently, $v=\rho(q-\log(\mathbf{s} h/\rho))$.
This log-domain formulation supports scalar evaluation without explicitly constructing $e^{u/\rho}$.

\subsection{Fixed-Metric Convergence}
\label{app:pduot-global}

We establish convergence guarantees for the fixed-metric core of \subsolver\ by verifying the error-bound conditions of restricted KL-UOT and combining them with existing Halpern and restart theory~\citep{yuan2020discerning,wittmann1992approximation,lieder2021convergence,lu2407restarted}.
The argument below proves Proposition~\ref{prop:pduot-convergence}; local sharpness is the intermediate property used to obtain the restart rate.
Fix $\mathcal N$, $\mathbf{D}_x$, and $\mathbf{D}_y$, and abbreviate $\mathbf{c}=\mathbf{c}_{\mathcal N}$ within this appendix.
All proximal operations in the proofs are exact.
For a positive definite matrix $\mathbf{P}$, write $\|\mathbf{z}\|_\mathbf{P}=(\mathbf{z}^\top \mathbf{P}\mathbf{z})^{1/2}$ and $\operatorname{dist}_\mathbf{P}(\mathbf{z},S)=\inf_{u\in S}\|\mathbf{z}-\mathbf{u}\|_\mathbf{P}$.
Define the restricted KKT mapping
\begin{equation}
\label{eq:pduot-kkt-map}
\mathcal M(\mathbf{x},\mathbf{y})=\begin{pmatrix}\mathbf{c}+\mathbf{A}_{\mathcal N}^\top \mathbf{y}+N_{\mathbb R_+^{|\mathcal N|}}(\mathbf{x})\\ \nabla G^*(\mathbf{y})-\mathbf{A}_{\mathcal N}\mathbf{x}\end{pmatrix},
\end{equation}
where $N_{\mathbb R_+^{|\mathcal N|}}$ is the convex normal-cone mapping.
The zeros of $\mathcal M$ are precisely the restricted KKT points, with $\mathbf{y}=-(\mathbf{f},\mathbf{g})$.
The mapping $\mathcal M$ is maximally monotone: $\partial F\times\partial G^*$ is maximally monotone, and the remaining linear map $(\mathbf{x},\mathbf{y})\mapsto(\mathbf{A}_{\mathcal N}^\top \mathbf{y},-\mathbf{A}_{\mathcal N}\mathbf{x})$ is skew-symmetric and defined everywhere.

For fixed positive diagonal steps, define the PDHG metric
\begin{equation*}
\mathbf{P}=\begin{pmatrix}\mathbf{D}_x^{-1}&-\mathbf{A}_{\mathcal N}^\top\\-\mathbf{A}_{\mathcal N}&\mathbf{D}_y^{-1}\end{pmatrix}.
\end{equation*}

\begin{proof}[Proof of the unrestarted statement in Proposition~\ref{prop:pduot-convergence}]
The restricted primal objective is continuous on the nonnegative orthant and coercive: both transported marginal totals equal $\|\mathbf x\|_1$, so the KL penalties grow superlinearly in $\|\mathbf x\|_1$, whereas the finite transport costs contribute a linear term.
It therefore attains a minimum $\mathbf x^\star$.
Node coverage makes its transported marginals strictly positive, as shown in Appendix~\ref{app:pduot-sharpness}.
Setting $\mathbf y^\star=\nabla G(\mathbf A_{\mathcal N}\mathbf x^\star)$, the primal first-order optimality condition and conjugacy give $0\in\mathcal M(\mathbf x^\star,\mathbf y^\star)$, so the restricted KKT solution set is nonempty.
The step-size condition makes $\mathbf P$ positive definite by the Schur complement.
For $\widehat{\mathbf{z}}=T\mathbf{z}$, the two proximal optimality conditions in~\eqref{eq:pduot-map} are
\begin{equation*}
\mathbf{D}_x^{-1}(\mathbf{x}-\widehat{\mathbf{x}})-\mathbf{A}_{\mathcal N}^\top \mathbf{y}\in\partial F(\widehat{\mathbf{x}}),\qquad
\mathbf{D}_y^{-1}(\mathbf{y}-\widehat{\mathbf{y}})+\mathbf{A}_{\mathcal N}(2\widehat{\mathbf{x}}-\mathbf{x})=\nabla G^*(\widehat{\mathbf{y}}).
\end{equation*}
Rearranging gives
\begin{equation}
\label{eq:pduot-resolvent}
\mathbf{P}(\mathbf{z}-T\mathbf{z})\in\mathcal M(T\mathbf{z}),\qquad T=(\mathbf{P}+\mathcal M)^{-1}\mathbf{P}.
\end{equation}
Monotonicity applied to the images of two points $\mathbf{z},\mathbf{u}$ yields
\begin{equation*}
\|T\mathbf{z}-T\mathbf{u}\|_\mathbf{P}^2\leq\langle T\mathbf{z}-T\mathbf{u},\mathbf{z}-\mathbf{u}\rangle_\mathbf{P}.
\end{equation*}
Thus $T$ is firmly nonexpansive, $R:=2T-\mathbf{I}$ is nonexpansive, and $\operatorname{Fix}R=\operatorname{Fix}T=\mathcal M^{-1}(0)$.
Equation~\eqref{eq:halpern}, initialized at $\mathbf{z}^0=\mathbf{z}^{\mathrm a}$, is the Halpern iteration for $R$ with coefficients $1/(k+2)$.
The Halpern convergence theorem~\citep{wittmann1992approximation} and the residual estimate of \citet[Theorem~2.1]{lieder2021convergence}, applied to $R$ in the $\mathbf P$ inner product, give convergence to a fixed point and
\begin{equation*}
\|\mathbf{z}^k-R\mathbf{z}^k\|_\mathbf{P}\leq\frac{2\operatorname{dist}_\mathbf{P}(\mathbf{z}^{\mathrm a},\operatorname{Fix}R)}{k+1}.
\end{equation*}
Since $\mathbf{z}^k-R\mathbf{z}^k=2(\mathbf{z}^k-T\mathbf{z}^k)$, we obtain
\begin{equation}
\label{eq:pduot-halpern-bound}
\|\mathbf z^k-T\mathbf z^k\|_{\mathbf P}\leq\frac{\operatorname{dist}_{\mathbf P}(\mathbf z^{\mathrm a},\operatorname{Fix}T)}{k+1}.
\end{equation}
Continuity of $T$ also gives convergence of the mapped points to the same KKT point.
\end{proof}

The degree-based steps in~\eqref{eq:pduot-metric} satisfy the required norm condition without estimating a spectral norm.
Indeed, Cauchy--Schwarz at each node gives, for any edge vector $\mathbf{h}$,
\begin{equation*}
\|\overline{\mathbf{D}}_y^{1/2}\mathbf{A}_{\mathcal N}\overline{\mathbf{D}}_x^{1/2}\mathbf{h}\|_2^2
=\frac{\zeta^2}{2}\sum_v\frac{1}{d_v}\left(\sum_{e\ni v}\mathbf{h}_e\right)^2
\leq\frac{\zeta^2}{2}\sum_v\sum_{e\ni v}\mathbf{h}_e^2
=\zeta^2\|\mathbf{h}\|_2^2.
\end{equation*}
Each edge has two endpoints, and the reciprocal factors $w$ in $\mathbf{D}_x$ and $\mathbf{D}_y$ cancel in the scaled incidence matrix.
Hence its norm is at most $\zeta<1$ for every fixed $w>0$.

\subsection{Positive Marginals and Local Sharpness}
\label{app:pduot-sharpness}

We first justify positivity of the transported marginals under the node-coverage assumption in Proposition~\ref{prop:pduot-convergence}.
Let $\mathbf{x}^\star$ be a restricted optimum and suppose its source marginal $r_i^\star$ is zero.
Choose an incident active edge $e=(i,j)$ and add flow $t>0$ on that edge.
The source KL change divided by $t$ is $\rho_a(\log(t/a_i)-1)$, which tends to $-\infty$ as $t\downarrow0$.
The cost quotient is finite, while the target KL quotient tends either to the finite value $\rho_b\log(q_j^\star/b_j)$ if $q_j^\star>0$, or to $-\infty$ otherwise.
The total quotient is therefore negative for sufficiently small $t$, contradicting optimality.
The target argument is symmetric, so $\mathbf{r}^\star,\mathbf{q}^\star>0$.

\paragraph{KKT error bound.}
Introduce the transported marginal vector $\mathbf m=(\mathbf r,\mathbf q)$ and write the restricted problem as
\begin{equation}
\label{eq:pduot-lifted-primal}
\min_{\mathbf m,\mathbf x}\;G(\mathbf m)+F(\mathbf x)\quad\text{s.t.}\quad\mathbf A_{\mathcal N}\mathbf x-\mathbf m=0.
\end{equation}
This problem satisfies the structured polyhedricity assumption of \citet[Assumptions~1.2--1.3]{yuan2020discerning}.
Indeed, the extended-valued KL function $G$ is proper, closed, and convex, is differentiable on the positive orthant, and has
\begin{equation*}
\nabla^2G(\mathbf m)=\operatorname{diag}(\rho_a/\mathbf r,\rho_b/\mathbf q)\succ0.
\end{equation*}
Thus $G$ is strongly convex on every compact convex subset of its subdifferential domain, its gradient is locally Lipschitz on the positive orthant, and the gradient norm diverges when approaching a finite boundary point.
In the notation of that assumption, take the smooth representation matrix to be the identity and the additional linear term to be zero.
The function $F$ is convex piecewise linear-quadratic, since it is a linear function plus the indicator of the nonnegative orthant.
Consequently, \citet[Theorem~59]{yuan2020discerning} gives metric subregularity of the lifted KKT mapping at every lifted KKT point, without rank or strict-complementarity assumptions.

To transfer this bound to~\eqref{eq:pduot-kkt-map}, define, for $\mathbf m>0$,
\begin{equation*}
\mathcal K(\mathbf m,\mathbf x,\mathbf y)=\begin{pmatrix}\nabla G(\mathbf m)-\mathbf y\\ \mathbf c+\mathbf A_{\mathcal N}^{\top}\mathbf y+N_{\mathbb R_+^{|\mathcal N|}}(\mathbf x)\\ \mathbf m-\mathbf A_{\mathcal N}\mathbf x\end{pmatrix}.
\end{equation*}
Up to signs and ordering of the residual blocks, this is the KKT mapping in the cited theorem.
Write $\mathcal S=\mathcal M^{-1}(0)$ and $\widetilde{\mathcal S}=\mathcal K^{-1}(0)$, and fix $\mathbf z^\star=(\mathbf x^\star,\mathbf y^\star)\in\mathcal S$.
The projection of $\widetilde{\mathcal S}$ onto $(\mathbf x,\mathbf y)$ is exactly $\mathcal S$.
For $\mathbf z=(\mathbf x,\mathbf y)$ near $\mathbf z^\star$, set $\mathbf m(\mathbf y)=\nabla G^*(\mathbf y)$.
This positive vector depends continuously on $\mathbf y$, and $\nabla G(\mathbf m(\mathbf y))=\mathbf y$.
The lifted error bound therefore yields, for some $\kappa>0$ and a neighborhood $\mathbf U$ of $\mathbf z^\star$,
\begin{equation}
\label{eq:pduot-kkt-error-bound}
\begin{aligned}
\operatorname{dist}_2(\mathbf z,\mathcal S)&\leq\operatorname{dist}_2((\mathbf m(\mathbf y),\mathbf x,\mathbf y),\widetilde{\mathcal S})\\
&\leq\kappa\operatorname{dist}_2(0,\mathcal K(\mathbf m(\mathbf y),\mathbf x,\mathbf y))\\
&=\kappa\operatorname{dist}_2(0,\mathcal M(\mathbf z)),\qquad \mathbf z\in\mathbf U.
\end{aligned}
\end{equation}
This bound measures distance to the entire solution set and allows nonunique optimal couplings, cycles in the positive-flow support, and edges with both zero flow and zero slack.

\paragraph{Fixed-point sharpness.}
Let $p_{\max}$ be the largest eigenvalue of $\mathbf P$.
For $\mathbf z$ sufficiently near $\mathbf z^\star$, continuity gives $T\mathbf z\in\mathbf U$.
Combining~\eqref{eq:pduot-kkt-error-bound} with~\eqref{eq:pduot-resolvent} and $\|\mathbf P\mathbf h\|_2\leq\sqrt{p_{\max}}\|\mathbf h\|_{\mathbf P}$ gives
\begin{equation*}
\operatorname{dist}_{\mathbf P}(T\mathbf z,\mathcal S)\leq\sqrt{p_{\max}}\kappa\|\mathbf P(\mathbf z-T\mathbf z)\|_2\leq\kappa p_{\max}\|\mathbf z-T\mathbf z\|_{\mathbf P}.
\end{equation*}
The triangle inequality then yields the local sharpness estimate
\begin{equation}
\label{eq:pduot-local-sharpness}
\alpha\operatorname{dist}_{\mathbf P}(\mathbf z,\mathcal S)\leq\|\mathbf z-T\mathbf z\|_{\mathbf P},\qquad \alpha=(1+\kappa p_{\max})^{-1}>0.
\end{equation}

\subsection{Local Linear Convergence with Restart}
\label{app:pduot-restart}

The following argument completes the restart statement in Proposition~\ref{prop:pduot-convergence}.
It combines~\eqref{eq:pduot-local-sharpness} with the sharpness--restart argument of \citet[Theorems~1--2 and Corollary~1]{lu2407restarted}, localized to a neighborhood of a restricted KKT point.
It uses the same fixed $\mathbf{P}$ across all epochs.
Write $\mathcal S=\operatorname{Fix}T$ and $\mathcal R(\mathbf{z})=\|\mathbf{z}-T\mathbf{z}\|_\mathbf{P}$.
Choose a closed $\mathbf{P}$-ball centered at $\mathbf{z}^\star$ contained in the sharpness neighborhood, and initialize the first anchor inside that ball.
Both $T$ and $2T-\mathbf{I}$ fix $\mathbf{z}^\star$ and are nonexpansive, so their images and the Halpern convex combinations remain in the ball.
Restarting at a mapped point preserves this property at every epoch.

For any epoch anchor $\mathbf{z}^{\mathrm a}$, equation~\eqref{eq:pduot-halpern-bound} and nonexpansiveness of $T$ give
\begin{equation}
\label{eq:pduot-epoch-bound}
\mathcal R(\mathbf{z}^k)\leq\frac{\operatorname{dist}_\mathbf{P}(\mathbf{z}^{\mathrm a},\mathcal S)}{k+1},\qquad
\mathcal R(T\mathbf{z}^k)\leq\mathcal R(\mathbf{z}^k).
\end{equation}
Applying sharpness at the next anchor $T\mathbf{z}^k$ yields
\begin{equation*}
\operatorname{dist}_\mathbf{P}(T\mathbf{z}^k,\mathcal S)\leq\frac{\operatorname{dist}_\mathbf{P}(\mathbf{z}^{\mathrm a},\mathcal S)}{\alpha(k+1)}.
\end{equation*}
For a fixed period of $L$ PDHG evaluations, take the next anchor to be $T\mathbf{z}^{L-1}$.
Any integer $L\geq e/\alpha$ therefore gives
\begin{equation*}
\operatorname{dist}_\mathbf{P}(\mathbf{z}^{\mathrm a}_{\ell+1},\mathcal S)\leq e^{-1}\operatorname{dist}_\mathbf{P}(\mathbf{z}^{\mathrm a}_{\ell},\mathcal S),
\end{equation*}
where $\ell$ indexes epochs.

Alternatively, fix $\vartheta\in(0,1)$ and restart at $T\mathbf{z}^k$ at the first index satisfying
\begin{equation}
\label{eq:pduot-proved-restart}
\mathcal R(\mathbf{z}^k)\leq\vartheta\mathcal R(\mathbf{z}^{\mathrm a}_{\ell}).
\end{equation}
If the anchor residual is zero, the anchor is already a KKT point.
Otherwise, sharpness at the anchor and~\eqref{eq:pduot-epoch-bound} show that this test passes once $k+1\geq1/(\alpha\vartheta)$.
Every epoch thus has at most $\lceil1/(\alpha\vartheta)\rceil$ PDHG evaluations, and
\begin{equation*}
\mathcal R(\mathbf{z}^{\mathrm a}_{\ell})\leq\vartheta^\ell\mathcal R(\mathbf{z}^{\mathrm a}_0),\qquad
\operatorname{dist}_\mathbf{P}(\mathbf{z}^{\mathrm a}_{\ell},\mathcal S)\leq\alpha^{-1}\vartheta^\ell\mathcal R(\mathbf{z}^{\mathrm a}_0).
\end{equation*}
These estimates prove linear convergence across restart epochs and, with the uniformly bounded epoch lengths, a geometric bound in total PDHG evaluations for the restart anchors.

The proofs establish a local rate once an anchor lies in the stated neighborhood; they do not assert that every practical restarted run enters it.
Changing $w$ changes $\mathbf{P}$, and the practical restart tests need not enforce~\eqref{eq:pduot-proved-restart}.
The full controller and finite-precision proximal evaluations are therefore outside these fixed-metric, exact-proximal guarantees.

\subsection{Practical Control Steps}
\label{app:pduot-controls}
\checkrestart tests fixed-point residual reduction relative to the epoch anchor, deterioration relative to the previous check after moderate reduction, and the epoch-length threshold described in the main text.
The practical reduction factors are $0.2$ and $0.5$, and the epoch-length threshold is $0.36$ times the total iteration count; the first restart check also starts a new epoch.
At a restart, let $d_x=\|\widehat{\mathbf{x}}-\mathbf{x}^{\mathrm a}\|_{\overline{\mathbf{D}}_x^{-1}}$ and $d_y=\|\widehat{\mathbf{y}}-\mathbf{y}^{\mathrm a}\|_{\overline{\mathbf{D}}_y^{-1}}$.
For valid, nonnegligible displacements, \updateweight forms $e=\log(d_x/d_y)-\log w$, updates an integral state $I\leftarrow\nu I+e$, and sets
\begin{equation*}
w^+ = \operatorname{clip}\bigl(\exp(\log w+k_pe+k_iI),w_{\min},w_{\max}\bigr).
\end{equation*}
When these displacement or residual checks are invalid, the implementation falls back to a previously retained weight with better primal--dual residual balance and resets the integral state.
Changing the metric starts a new epoch; the implementation refreshes the mapped point and residual baseline under that metric.
Algorithm~\ref{alg:pduot} presents the control structure with one outer iteration per restart epoch.
For clarity, its stopping and restart tests are shown at each iteration; the implementation evaluates them periodically, reusing the last residual between checks.

\section{Stopping Criterion}
\label{app:stopping-criterion}

This section defines the three measures used in~\eqref{eq:pduot-stopping}, then analyzes the gap and its balanced-OT limit.
All stopping measures are evaluated in the original problem coordinates on a fixed active support $\mathcal N$.

\subsection{Residuals and Termination}

For a PDHG input $\mathbf{z}=(\mathbf{x},\mathbf{y})$ and its mapped point $\widehat{\mathbf{z}}=(\widehat{\mathbf{x}},\widehat{\mathbf{y}})=T\mathbf{z}$, write $(\mathbf{r},\mathbf{q})=\mathbf{A}_{\mathcal N}\widehat{\mathbf{x}}$ and $(\mathbf{f},\mathbf{g})=-\widehat{\mathbf{y}}$.
The marginal responses are $\widehat{\mathbf{r}}=\mathbf{a}\odot e^{-\mathbf{f}/\rho_a}$ and $\widehat{\mathbf{q}}=\mathbf{b}\odot e^{-\mathbf{g}/\rho_b}$, and the restricted dual slack is $\mathbf{s}_{\mathcal N}=\mathbf{c}_{\mathcal N}-\mathbf{A}_{\mathcal N}^\top(\mathbf{f},\mathbf{g})$.
We define the marginal-consistency and stationarity residuals by
\begin{equation}
\label{eq:pduot-residuals}
\begin{aligned}
\delta_{\mathrm m}&=\frac{\|(\mathbf{r}-\widehat{\mathbf{r}},\mathbf{q}-\widehat{\mathbf{q}})\|}{1+\|(\mathbf{a},\mathbf{b})\|},\\
\delta_{\mathrm s}&=\frac{\|\mathbf{s}_{\mathcal N}-\boldsymbol{\lambda}\|}{1+\|\mathbf{c}_{\mathcal N}\|},
\end{aligned}
\end{equation}
where $\boldsymbol{\lambda}=\mathbf{D}_x^{-1}(\widehat{\mathbf{x}}-\mathbf{x})+\mathbf{c}_{\mathcal N}+\mathbf{A}_{\mathcal N}^\top \mathbf{y}$ is the bound multiplier supplied by the primal projection.
The projection ensures $\widehat{\mathbf{x}}\geq0$, $\boldsymbol{\lambda}\geq0$, and $\lambda_e\widehat x_e=0$.
Consequently, zero stationarity residual implies restricted dual feasibility and complementarity, while zero marginal-consistency residual gives the KL marginal optimality relations.
Together these conditions imply the restricted KKT conditions.

Let $t=\langle \mathbf{c}_{\mathcal N},\widehat{\mathbf{x}}\rangle$ and $d=\langle \mathbf{f},\widehat{\mathbf{r}}\rangle+\langle \mathbf{g},\widehat{\mathbf{q}}\rangle$.
The relative gap is
\begin{equation}
\label{eq:pduot-split-gap}
\delta_{\mathrm g}=\frac{|t-d|}{1+|t|+|d|}.
\end{equation}
\checkstop returns the mapped point only when all three measures satisfy~\eqref{eq:pduot-stopping}.
The residual norms are $\ell_2$ by default, with an $\ell_\infty$ option at the finest hierarchy level.
These tests concern the restricted problem; Algorithm~\ref{alg:ours} separately checks dual violations outside the active support through full-space pricing.

\subsection{Gap Criterion and the Balanced Limit}
\label{app:gap-criterion}

For the following analysis, $\mathbf{x}$ denotes a generic nonnegative evaluation point, including the mapped coupling used above.
Let $\Phi_{\mathcal N}(\mathbf{x})$ be the restricted primal objective, and let $D(\mathbf{f},\mathbf{g})$ be the dual objective in~\eqref{eq:dual-uot} evaluated at the current potentials.
Write $(\mathbf{r},\mathbf{q})=\mathbf{A}_{\mathcal N}\mathbf{x}$, $\widehat{\mathbf{r}}=\mathbf{a}\odot e^{-\mathbf{f}/\rho_a}$, and $\widehat{\mathbf{q}}=\mathbf{b}\odot e^{-\mathbf{g}/\rho_b}$.

\paragraph{Original objective gap.}
The generalized-KL identity
\begin{equation}
\rho_a\KL(\mathbf{r}\,|\,\mathbf{a})+\rho_a\sum_i a_i(e^{-f_i/\rho_a}-1)
=-\langle \mathbf{f},\mathbf{r}\rangle+\rho_a\KL(\mathbf{r}\,|\,\widehat{\mathbf{r}})
\end{equation}
and its target counterpart give the exact decomposition
\begin{equation}
\label{eq:original-gap-decomposition}
\Phi_{\mathcal N}(\mathbf{x})-D(\mathbf{f},\mathbf{g})
=\langle \mathbf{c}_{\mathcal N},\mathbf{x}\rangle-\langle \mathbf{f},\mathbf{r}\rangle-\langle \mathbf{g},\mathbf{q}\rangle
+\rho_a\KL(\mathbf{r}\,|\,\widehat{\mathbf{r}})+\rho_b\KL(\mathbf{q}\,|\,\widehat{\mathbf{q}}).
\end{equation}
When the potentials are dual feasible on $\mathcal N$, the first three terms form a nonnegative complementarity term, and the entire expression is a primal--dual gap bounding restricted objective suboptimality.
At arbitrary infeasible dual iterates, it remains an objective difference without this bound.
The divergence terms weight marginal-response errors by the penalty coefficients.
For positive $\widehat{\mathbf{r}}$ and $\mathbf{r}$ near $\widehat{\mathbf{r}}$, for example,
\begin{equation}
\rho_a\KL(\mathbf{r}\,|\,\widehat{\mathbf{r}})
=\frac{\rho_a}{2}\sum_i\frac{(r_i-\widehat r_i)^2}{\widehat r_i}
+o(\|\mathbf{r}-\widehat{\mathbf{r}}\|_2^2),
\end{equation}
where the expansion holds for fixed $\rho_a$ and $\widehat{\mathbf{r}}$.
Thus a prescribed marginal residual tolerance does not give a penalty-independent bound on the original objective gap.

\paragraph{Gap for the split optimality system.}
Introducing independent transported marginals yields the equivalent split formulation
\begin{equation}
\min_{\mathbf{x},\mathbf{r},\mathbf{q}\geq0}\;\langle \mathbf{c}_{\mathcal N},\mathbf{x}\rangle+\rho_a\KL(\mathbf{r}\,|\,\mathbf{a})+\rho_b\KL(\mathbf{q}\,|\,\mathbf{b})
\quad\text{s.t.}\quad \mathbf{A}_{\mathcal N}\mathbf{x}=(\mathbf{r},\mathbf{q}).
\end{equation}
Its marginal responses are the minimizers of the marginal Lagrangian terms for given $(\mathbf{f},\mathbf{g})$.
At these responses, $\nabla(\rho_a\KL)(\widehat{\mathbf{r}}\,|\,\mathbf{a})=-\mathbf{f}$ and $\nabla(\rho_b\KL)(\widehat{\mathbf{q}}\,|\,\mathbf{b})=-\mathbf{g}$.
Pairing the stationarity relations with primal variables motivates the signed consistency quantity
\begin{equation}
H=\langle \mathbf{c}_{\mathcal N},\mathbf{x}\rangle-\langle \mathbf{f},\widehat{\mathbf{r}}\rangle-\langle \mathbf{g},\widehat{\mathbf{q}}\rangle,
\qquad
\delta_{\mathrm g}=\frac{|H|}{1+|\langle \mathbf{c}_{\mathcal N},\mathbf{x}\rangle|+|\langle \mathbf{f},\widehat{\mathbf{r}}\rangle+\langle \mathbf{g},\widehat{\mathbf{q}}\rangle|}.
\end{equation}
This is the gap in~\eqref{eq:pduot-split-gap}, expressed at a generic evaluation point.
Its exact relation to the original objective difference is
\begin{equation}
\Phi_{\mathcal N}(\mathbf{x})-D(\mathbf{f},\mathbf{g})
=H-\langle \mathbf{f},\mathbf{r}-\widehat{\mathbf{r}}\rangle-\langle \mathbf{g},\mathbf{q}-\widehat{\mathbf{q}}\rangle
+\rho_a\KL(\mathbf{r}\,|\,\widehat{\mathbf{r}})+\rho_b\KL(\mathbf{q}\,|\,\widehat{\mathbf{q}}).
\end{equation}
At exact marginal consistency, $H$ equals the original objective difference.
At finite tolerance, $\delta_{\mathrm g}$ measures consistency of the split system and is combined with $\delta_{\mathrm m}$ and $\delta_{\mathrm s}$.
In particular, exact marginal consistency and stationarity with a nonnegative complementary bound multiplier imply the restricted KKT conditions and global restricted optimality.
The scalar $\delta_{\mathrm g}$ alone does not bound original-objective suboptimality.

\paragraph{Compatibility with balanced OT.}
Consider equal-total-mass references and a support admitting a balanced transport plan.
For fixed bounded potentials, as $\rho_a,\rho_b\to\infty$ the marginal responses approach $\mathbf{a}$ and $\mathbf{b}$, respectively.
Consequently, the gap expression tends to
\begin{equation}
\frac{|\langle \mathbf{c}_{\mathcal N},\mathbf{x}\rangle-\langle \mathbf{a},\mathbf{f}\rangle-\langle \mathbf{b},\mathbf{g}\rangle|}
{1+|\langle \mathbf{c}_{\mathcal N},\mathbf{x}\rangle|+|\langle \mathbf{a},\mathbf{f}\rangle+\langle \mathbf{b},\mathbf{g}\rangle|},
\end{equation}
and the marginal-response residual becomes the balanced marginal-feasibility residual.
These are the corresponding balanced-OT stopping quantities.
The original dual objective also has the balanced limit at bounded potentials; the distinction is the nonuniform penalty weighting of marginal errors in~\eqref{eq:original-gap-decomposition} along approximately feasible primal iterates.

For a concrete example, take one edge with $a=b=1$, $c=f=g=0$, $x=1+\delta>0$, and equal penalties $\rho$.
Then $\widehat r=\widehat q=1$, $\delta_{\mathrm g}=\delta_{\mathrm s}=0$, and $\delta_{\mathrm m}=\sqrt{2}|\delta|/(1+\sqrt{2})$ under the $\ell_2$ policy.
Meanwhile,
\begin{equation}
\Phi_{\mathcal N}(x)-D(f,g)=2\rho\bigl((1+\delta)\log(1+\delta)-\delta\bigr),
\end{equation}
which is asymptotic to $\rho\delta^2$ for small $\delta$.
For any fixed nonzero $\delta$, the usual relative original-objective gap tends to one as $\rho\to\infty$, even though the marginal residual remains at a fixed small scale.
Thus imposing a small original-objective gap can require increasingly accurate marginals as the penalties grow.
\subsolver instead applies a common tolerance to the three split optimality measures, preserving the balanced-OT stopping scale; original-objective gaps are evaluated separately when required for numerical comparisons.

\section{Natural Extension to UOT with general Marginal Penaltys}
\label{app:extension_penalty}
% CN: 本附录预留 squared-L2 边际惩罚版本的定义、算法适配和理论说明。
% EN: This appendix is reserved for the squared-L2 marginal-penalty variant, its algorithmic adaptation, and supporting analysis.

\section{Sparse Entropic Refinement}
\label{app:entropic-extension}

On a fixed support $\mathcal{N}$, define $K_{ij}=a_i b_j\exp(-C_{ij}/\varepsilon)$ for $(i,j)\in\mathcal{N}$ and zero otherwise.
The supported coupling has the form $X_{ij}=u_iK_{ij}v_j$, with generalized scaling updates
\begin{equation}
\mathbf{u}\leftarrow\left(\frac{\mathbf{a}}{\mathbf{K}\mathbf{v}}\right)^{\rho_a/(\rho_a+\varepsilon)},\qquad
\mathbf{v}\leftarrow\left(\frac{\mathbf{b}}{\mathbf{K}^\top \mathbf{u}}\right)^{\rho_b/(\rho_b+\varepsilon)}.
\end{equation}
Our implementation evaluates the scaling iterations in the log domain using sparse row and column reductions on the GPU.
The KL term in~\eqref{eq:sparse-entropic-extension} uses the full reference $\mathbf{a}\otimes \mathbf{b}$, including its constant contribution outside $\mathcal{N}$, when evaluating the objective.

The default choice is $\mathcal{N}=\mathcal{N}_\star$.
An optional expansion uses the positive-flow support of the unregularized solution $\mathbf{X}^0$ and bidirectional top-$b_{\mathrm e}$ selection under its final dual scores:
\begin{equation}
\mathcal{N}_{b_{\mathrm e}}=\supp(\mathbf{X}^0)\cup\operatorname{RowTop}_{b_{\mathrm e}}(\boldsymbol{\sigma})\cup\operatorname{ColTop}_{b_{\mathrm e}}(\boldsymbol{\sigma}).
\end{equation}
Here, $b_{\mathrm e}$ denotes the per-node edge budget.
Constructing this alternative support requires an additional full-space scan.

\section{Two-Coupling UFGW Updates}
\label{app:ufgw-extension}

\paragraph{Reduction to KL-UOT.}
Write $|\mathbf{Q}|=\sum_{kl}Q_{kl}$ and $\mathbf{q}_a=\mathbf{Q}\mathbf{1}$, $\mathbf{q}_b=\mathbf{Q}^\top\mathbf{1}$.
The tensorized marginal term satisfies
\begin{equation}
\KL(\mathbf{p}\otimes \mathbf{q}\,|\,\mathbf{a}\otimes \mathbf{a})
=|\mathbf{q}|\KL(\mathbf{p}\,|\,\mathbf{a})+|\mathbf{p}|\sum_kq_k\log(q_k/a_k)+|\mathbf{a}|(|\mathbf{a}|-|\mathbf{q}|).
\end{equation}
Consequently, fixing $\mathbf{Q}$ in~\eqref{eq:ufgw-relaxation} yields, up to constants independent of $\mathbf{P}$,
\begin{equation}
\label{eq:ufgw-block}
\min_{P\geq0}\;\langle \mathbf{C}^\mathbf{Q},\mathbf{P}\rangle+\rho_a|\mathbf{Q}|\KL(\mathbf{P}\mathbf{1}\,|\,\mathbf{a})+\rho_b|\mathbf{Q}|\KL(\mathbf{P}^\top\mathbf{1}\,|\,\mathbf{b}),
\end{equation}
where
\begin{equation}
\begin{aligned}
\mathbf{C}^Q_{ij}&=(1-\alpha)\sum_{kl}(A_{ik}-B_{jl})^2Q_{kl}+\frac{\alpha}{2}M_{ij}+\kappa_Q,\\
\kappa_Q&=\rho_a\sum_k(\mathbf{q}_a)_k\log((\mathbf{q}_a)_k/a_k)+\rho_b\sum_l(\mathbf{q}_b)_l\log((\mathbf{q}_b)_l/b_l).
\end{aligned}
\end{equation}
The $\mathbf{Q}$ block is obtained symmetrically; our structural matrices are symmetric squared-Euclidean distances.
The prototype initializes with the product coupling $\mathbf{a}\otimes \mathbf{b}$, stored implicitly, and alternates the two block solves.
After solving the $\mathbf{P}$ block, it applies $\mathbf{P}\leftarrow\sqrt{|\mathbf{Q}|/|\mathbf{P}|}\,\mathbf{P}$; after solving the $\mathbf{Q}$ block, it applies the symmetric rescaling using the current $\mathbf{P}$.
These steps match the block-coordinate implementation used in the experiments.
The outer procedure optimizes a nonconvex two-coupling relaxation; the restricted-solver convergence results do not imply global optimality of this relaxation or equality of the two returned couplings.

\paragraph{Entropic blocks.}
For normalized reference masses, the optional entropic objective adds $\varepsilon\KL(\mathbf{P}\otimes \mathbf{Q}\,|\,(\mathbf{a}\otimes \mathbf{b})\otimes(\mathbf{a}\otimes \mathbf{b}))$ to~\eqref{eq:ufgw-relaxation}.
With $\mathbf{Q}$ fixed, the block therefore acquires $\varepsilon|\mathbf{Q}|\KL(\mathbf{P}\,|\,\mathbf{a}\otimes \mathbf{b})$ and the scalar cost contribution $\varepsilon\sum_{kl}Q_{kl}\log(Q_{kl}/(a_kb_l))$.
We solve its unregularized counterpart with \ours and apply sparse entropic refinement on the resulting support.
The scalar cost translation can be absorbed into rescaled KL reference masses, retaining the same block minimizer without modifying the squared-distance solver.

\paragraph{Implicit structural costs.}
For squared-Euclidean structural matrices, write exact factorizations $\mathbf{A}=\mathbf{U}_A\mathbf{V}_A^\top$ and $\mathbf{B}=\mathbf{U}_B\mathbf{V}_B^\top$.
The linearized structural cost is
\begin{equation}
L_Q=(\mathbf{A}\odot \mathbf{A})\mathbf{q}_a\mathbf{1}_n^\top+\mathbf{1}_m((\mathbf{B}\odot \mathbf{B})\mathbf{q}_b)^\top-2\mathbf{A}\mathbf{Q}\mathbf{B}^\top.
\end{equation}
Its cross term depends on the small moment matrix $\mathbf{V}_A^\top \mathbf{Q}\mathbf{V}_B$, and the two squared-distance terms can likewise be evaluated through small weighted moments.
More explicitly, the matrix
\begin{equation}
\mathbf{H}_Q=\begin{pmatrix}
\mathbf{V}_A^\top\operatorname{diag}(\mathbf{q}_a)\mathbf{V}_A&\mathbf{V}_A^\top \mathbf{Q}\mathbf{V}_B\\
\mathbf{V}_B^\top \mathbf{Q}^\top \mathbf{V}_A&\mathbf{V}_B^\top\operatorname{diag}(\mathbf{q}_b)\mathbf{V}_B
\end{pmatrix}
\end{equation}
is positive semidefinite because it is a sum of outer products weighted by $Q_{kl}$.
Factoring $\mathbf{H}_Q=\mathbf{F}\mathbf{F}^\top$ and setting $z_i=(\mathbf{U}_{A,i},0)\mathbf{F}$, $\mathbf{w}_j=(0,\mathbf{U}_{B,j})\mathbf{F}$ gives $(L_Q)_{ij}=\|z_i-\mathbf{w}_j\|_2^2$.
Scaling these embeddings by $\sqrt{1-\alpha}$ and concatenating feature vectors scaled by $\sqrt{\alpha/2}$ represents the nonconstant block cost exactly.
For fixed structural and feature dimensions, the embedding storage grows linearly with the number of points and the moment accumulation depends on the maintained coupling support.

\end{document}
